% ======================================================================
% PUV Pipeline -- User Guide and Reference
% Single, self-contained writeup. Compile with:
%     pdflatex PUV_Pipeline_Guide.tex   (run twice for the table of contents)
% No external files or bibliography needed.
% ======================================================================
\documentclass[11pt]{article}
\usepackage[margin=1in]{geometry}
\usepackage{amsmath}
\usepackage{booktabs}
\usepackage{array}
\usepackage{xcolor}
\usepackage{listings}
\usepackage[colorlinks=true,linkcolor=blue!50!black,urlcolor=blue!50!black]{hyperref}
\usepackage[utf8]{inputenc}

\definecolor{codebg}{gray}{0.96}
\definecolor{codekw}{rgb}{0.35,0.0,0.55}
\definecolor{codecm}{rgb}{0.30,0.50,0.30}
\lstset{
  basicstyle=\ttfamily\small,
  backgroundcolor=\color{codebg},
  keywordstyle=\color{codekw}\bfseries,
  commentstyle=\color{codecm}\itshape,
  frame=single, framesep=5pt, rulecolor=\color{gray!40},
  breaklines=true, showstringspaces=false,
  language=Matlab, columns=fullflexible, keepspaces=true,
  xleftmargin=6pt, xrightmargin=6pt,
}
\newcommand{\code}[1]{\texttt{#1}}
\newcommand{\qcflag}[1]{\texttt{qc\_flag}\,=\,#1}
\providecommand{\degree}{\ensuremath{^{\circ}}}

\title{\textbf{Nortek Vector PUV Pipeline}\\[2pt]
       \large User Guide and Reference}
\author{Holden Leslie-Bole \\ Scripps Institution of Oceanography}
\date{July 2026}

\begin{document}
\maketitle

\begin{abstract}
\noindent
This is the entry point for anyone using the Nortek Vector PUV pipeline. It takes raw
pressure--velocity (PUV) instrument files through five processing levels into
quality-controlled time series, wave spectra, and nearshore wave-forcing diagnostics, using
one deployment-agnostic configuration system. Sections \ref{sec:start}--\ref{sec:running}
are a step-by-step guide to running your own deployment. Section \ref{sec:qc} documents the
quality-control and data-provenance system, including the per-channel QC that recovers
storm-peak velocity from partial sensor failures. Sections \ref{sec:methods}--\ref{sec:outputs}
are the methods and output reference. You do not need to read it front to back: to just run
the code, read \S\ref{sec:start}--\ref{sec:running} and \S\ref{sec:qc}.
\end{abstract}

\tableofcontents
\newpage

% ======================================================================
\section{What the pipeline is}
\label{sec:intro}

Nearshore wave and current dynamics are measured with bottom-mounted Nortek Vector acoustic
Doppler velocimeters (ADVs) in pressure--velocity mode, sampling continuously at 2~Hz. The
pipeline turns the raw instrument files (\code{.dat}/\code{.sen}/\code{.hdr}) into standardized
outputs at five levels. Every deployment runs through the same drivers and produces the same
output structures, so results are comparable across sites and years.

\begin{center}
\begin{tabular}{@{}c l p{9.3cm}@{}}
\toprule
Level & Output & What it contains \\
\midrule
\textbf{L1} & QC'd 2~Hz time series & Burst merge, clock-drift and tilt correction,
  per-channel quality control, rotation to a geographic frame \\
\textbf{L2} & Per-segment spectra & Multi-taper power spectra, pressure correction to surface
  elevation, bulk wave parameters ($H_s$, $T_p$, direction), near-bed velocity, bed and
  Reynolds stress, velocity moments \\
\textbf{L3} & Forcing metrics & Frequency-band energy decomposition, storm detection,
  transport proxies (Shields, Rouse), tidal and undertow current decomposition \\
\textbf{L4} & Nonlinear / IG diagnostics & Surface-elevation bands, incident/reflected
  infragravity split, bispectra (skewness, asymmetry, bicoherence), velocity PDFs \\
\textbf{L5} & \emph{(planned)} & PUV--altimeter integration: bed change vs.\ forcing \\
\bottomrule
\end{tabular}
\end{center}

\noindent
L1--L3 are universal. L4 adds nonlinear-wave and infragravity diagnostics on top of L2. Each
level reads the previous level's \code{.mat} files, so you always run L1$\rightarrow$L2$\rightarrow$L3$\rightarrow$L4
in order.

% ======================================================================
\section{Getting started}
\label{sec:start}

\subsection{Requirements}

\begin{itemize}
\item \textbf{MATLAB} R2023b or newer, with the \textbf{Signal Processing Toolbox}
  (multi-taper and Welch spectra) and the \textbf{Mapping Toolbox}
  (\code{igrfmagm} for magnetic declination). \code{t\_tide} must be on the path for L3
  tidal analysis. The Aerospace Toolbox is \emph{not} required.
\item \textbf{Internet at runtime}: L2 fetches the shore-normal angle from the CDIP THREDDS
  server (it falls back to buoy coordinates if unavailable).
\item \textbf{Lab server access} for raw data and deployment metadata:
  \code{/Volumes/group/PUV\_data/Vector/} and \code{/Volumes/group/DeploymentNotes/}.
\end{itemize}

\subsection{First run}

\begin{lstlisting}
% 1. From the repo root, put every pipeline directory on the MATLAB path:
cd /path/to/PUV_Pipeline
startup_puv

% 2. (Recommended) cache the raw files locally -- network reads are slow, and the
%    drivers automatically use raw_cache/ when it exists:
reg = deployment_registry();
cfg = reg('TBR23')();
copy_raw_to_local(cfg);

% 3. Run the levels in order. Set the deployment name at the top of each driver:
PUV_L1_driver    % raw  -> QC'd timeseries      (outputs/L1/TBR23/)
PUV_L2_driver    % spectra + bulk params        (outputs/L2/TBR23/)
PUV_L3_driver    % forcing characterization     (outputs/L3/TBR23/)
PUV_L4_driver    % nonlinear / IG dynamics      (outputs/L4/TBR23/)
\end{lstlisting}

\noindent
To process every registered deployment, use the \code{*\_run\_all} batch scripts
(\code{PUV\_L1\_run\_all}, \code{PUV\_L2\_run\_all}, \dots), which skip deployments that
already have outputs. Outputs land in
\code{outputs/L\{1..4\}/\{deployment\}/\{label\}\_L\{n\}.mat} and are not version-controlled;
they are reproducible from raw data plus code.

% ======================================================================
\section{Configuring a deployment}
\label{sec:config}

Every deployment is registered in \code{config/deployment\_registry.m}, which maps a short
name to a config function --- for example \code{'TBR23'} $\rightarrow$ \code{@TBR23\_config}.
The naming convention is \code{SITE + year + optional season letter} (\code{TBR23} = Torrey
Pines summer 2023; \code{TOR24S} = Torrey spring 2024; \code{SIO24A--C} = SIO Pier 2024
chronological).

\subsection{Adding a new deployment}

\begin{enumerate}
\item Copy an existing \code{config/<SITE>\_config.m} as a template.
\item Fill in the processing parameters from the authoritative source,
  \code{/Volumes/group/DeploymentNotes/DeploymentNotes\{year\}.xls}: per-instrument
  \code{latlon}, \code{heading}, \code{clockDrift}, and \code{doffp} (pressure-sensor height
  above the bed). The annual checkout spreadsheets have serial numbers but not these
  numerical parameters.
\item Register the deployment in \code{config/deployment\_registry.m}.
\item Check \code{config/CONFIG\_REVIEW\_NOTES.md} and \code{config/DOFFP\_LOOKUP\_CHECKLIST.md}
  for deployment-specific gotchas.
\end{enumerate}

\subsection{Quality-control options: \texttt{cfg.qcOpts}}
\label{sec:qcopts}

The L1 quality control has one option struct, \code{cfg.qcOpts}, with sensible defaults. You
only need to set it deliberately in one case, but that case matters (see the boxed note).
The fields:

\begin{center}
\begin{tabular}{@{}l l p{7.7cm}@{}}
\toprule
Field & Default & Meaning \\
\midrule
\code{corrMin}    & 70 & Minimum beam correlation (\%), Nortek's recommendation. \\
\code{Tvalid}     & \code{[-2 40]} & \textbf{Plausible water-temperature range (\degree C).}
  The primary thermistor-failure discriminant. See the boxed note. \\
\code{maxJump}    & \code{Inf} & Optional per-sample temperature-step gate; \code{Inf} = off.
  Catches instantaneous glitches, never genuine drift. \\
\code{tiltStdMax} & 2  & Rolling-std tilt-variability threshold (\degree). \\
\code{tiltAbsMax} & 30 & Absolute tilt = toppled frame; always invalidates velocity. \\
\code{tiltWindow} & 120 & Samples for the rolling std (60~s at 2~Hz). \\
\code{cFactorTol} & 0.002 & Below this, no sound-speed rescale is applied. \\
\bottomrule
\end{tabular}
\end{center}

\begin{center}
\fbox{\parbox{0.94\textwidth}{\small
\textbf{Set \code{Tvalid} to the site's water-temperature range.} The default
\code{[-2 40]} is a wide safety bound that catches gross failures (freezing, firmware
garbage) anywhere on Earth, but it treats an implausibly cold in-water reading from a failing
thermistor as valid. For San Diego coastal deployments set
\code{cfg.qcOpts.Tvalid = [9 26]}; this catches the subtler within-bounds thermistor failures
(a December reading of $8\,\degree$C at Torrey Pines is not real water) without ever
mis-flagging genuine seasonal or upwelling temperatures. Section~\ref{sec:qc} explains why.
}}
\end{center}

% ======================================================================
\section{Running the pipeline, level by level}
\label{sec:running}

Each level has a driver (\code{PUV\_L\{n\}\_driver.m}) with a \code{deployment\_name} variable
at the top. Set it and run. The batch versions (\code{PUV\_L\{n\}\_run\_all.m}) loop over the
registry.

\begin{description}
\item[L1 --- raw to QC'd time series.] Reads the raw \code{.dat}/\code{.sen}/\code{.hdr},
  merges bursts to a continuous 2~Hz series, corrects clock drift, applies per-channel QC
  (\S\ref{sec:qc}), corrects tilt, and rotates velocity to the geographic buoy frame ($+x$
  west, $+y$ north, $+z$ up). Writes one \code{*\_processed.mat} per instrument holding a
  \code{PUV} struct with the time series and a \code{PUV.qc} provenance sub-struct.
\item[L2 --- spectra and wave statistics.] Segments the L1 series into 1-hour windows,
  estimates multi-taper spectra, converts pressure to surface elevation, and computes bulk
  wave parameters, near-bed velocity, bed and Reynolds stress, and velocity moments. Writes an
  \code{L2} struct, one row per segment. This is the level most analyses read.
\item[L3 --- forcing characterization.] Decomposes energy into frequency bands, detects storm
  events (with MOP gap-filling), computes transport proxies (Shields, Rouse), and separates
  tidal from subtidal (wave-driven) currents. Writes an \code{L3} struct.
\item[L4 --- nonlinear / IG dynamics.] Adds surface-elevation band products, an
  incident/reflected infragravity split, bispectra, and velocity PDFs on top of L2. Writes an
  \code{L4} struct.
\end{description}

\noindent
Each driver errors with a clear message if the previous level's output is missing, so run them
in sequence. Runtime is dominated by reading raw files over the network; cache locally first
(\code{copy\_raw\_to\_local}).

% ======================================================================
\section{Quality control and data provenance}
\label{sec:qc}

This is the part that most affects whether your results are trustworthy, so read it even if
you skip the methods. The guiding principle is \textbf{channel decoupling}: the pressure,
velocity, temperature, and tilt channels are judged independently, and one dead channel never
discards data from a channel that is still good.

\subsection{Why channel decoupling}

A Nortek Vector has a Doppler velocity sensor and an auxiliary sensor block (pressure,
thermistor, compass). These fail independently. The pipeline's original QC nulled an entire
data row whenever \emph{any} channel failed, on the assumption that a compromised instrument
gives unreliable velocity. That assumption is often false. At Torrey Pines in December 2023,
one frame's pressure sensor and thermistor died during a storm while its Doppler channel
stayed healthy (beam correlations 92--97\%); the old QC threw away the velocity anyway,
including every hour of the largest waves in the deployment. Roughly 120~hours of
storm-peak velocity were recoverable and had been discarded.

The current pipeline judges each channel on its own evidence and records what it did, so a
partial failure costs only the channel that actually failed.

\subsection{Per-channel validity}

L1 computes four per-sample validity masks, using the thresholds in \code{cfg.qcOpts}:

\begin{itemize}
\item \textbf{Velocity} is valid when the minimum beam correlation is $\geq$ \code{corrMin}
  and the frame is not toppled (see tilt, below). A dead pressure sensor or thermistor does
  \emph{not} invalidate velocity.
\item \textbf{Pressure} is valid inside $[\overline{P}/2,\, 2\overline{P}]$, where
  $\overline{P}$ is a reference median from a healthy in-water window.
\item \textbf{Temperature} is valid inside the physical range \code{Tvalid}, optionally with
  a per-sample step gate (\code{maxJump}).
\item \textbf{Tilt} is judged two ways. A \emph{variability} spike (rolling std above
  \code{tiltStdMax}) is a noisy reading; an \emph{absolute} tilt beyond \code{tiltAbsMax}
  ($30\degree$) means the frame has toppled or moved.
\end{itemize}

\noindent
\textbf{Tilt is not an auxiliary channel.} Pressure and temperature say nothing about the
Doppler measurement, so they cannot gate it. Tilt is different: it describes the measurement
geometry. A toppled frame contaminates the velocity with its own motion, and no rotation
repairs that, so an absolute tilt beyond $30\degree$ \emph{always} invalidates velocity ---
regardless of the state of any other channel. A variability spike is trusted only when the
sensor block it lives on is healthy; when the block has failed, the jittery tilt reading is
discarded and a static tilt from the healthy part of the record is substituted for the
rotation.

\subsection{The temperature discriminant}
\label{sec:tvalid}

The thermistor matters more than it looks, because the instrument uses the measured sound
speed --- which it computes from temperature --- to scale recorded velocity. A dead thermistor
therefore biases every velocity linearly. The discriminant for ``failed'' is
\emph{physical plausibility} (\code{Tvalid}), not a distance from the record mean. A
distance-from-mean test would flag a genuine cold-upwelling event as a failure and corrupt its
velocity; a physical-range test never does, because real seawater at a real temperature is
always valid. This is why \code{Tvalid} should be set to the site's water range
(\S\ref{sec:qcopts}): the narrower the physically-justified range, the more subtle failures it
catches, with no risk to good data.

\subsection{Sound-speed correction}

Where the thermistor has failed but the velocity is good, the velocity is rescaled to remove
the sound-speed bias, following the Nortek manual (N3015-030, \S2.4.9):
\[
  V_\text{corrected} = V_\text{recorded}\,\frac{c_\text{true}}{c_\text{measured}} .
\]
The reference sound speed $c_\text{true}$ comes from a companion frame if the config supplies
one (\code{cfg.soundSpeedRef}), otherwise from the instrument's own healthy
temperature--sound-speed relation, evaluated at a reference temperature and never
extrapolated beyond it. Where the thermistor is healthy the factor is exactly~1, so the
correction is a no-op on good data --- a property, not a threshold.

\subsection{Provenance flags}

Nothing reconstructed or rescaled is ever presented as a clean measurement. Every L2 segment
carries provenance so downstream analyses can filter:

\begin{center}
\begin{tabular}{@{}l p{11cm}@{}}
\toprule
Field & Meaning \\
\midrule
\code{segValid}      & Old-style ``all channels good'' flag; unchanged meaning, so existing
  analyses keyed on it behave exactly as before. \\
\code{segValid\_vel} & Velocity products (moments, mean flow, Reynolds stress) are usable. \\
\code{segValid\_p}   & Pressure products ($H_s$, $S_\eta$, depth) are usable. \\
\code{qc\_flag}      & QARTOD-style: 1~good, 2~not evaluated, 3~suspect (recovered or
  sound-speed-rescaled), 4~fail (physically implausible). Anything reconstructed is~3,
  never~1. \\
\code{Hs\_source}    & \code{'measured'} / \code{'reconstructed'} / \code{'none'}. \\
\code{vel\_c\_factor}, \code{vel\_c\_corrected} & The sound-speed factor applied (exactly~1 on
  healthy data) and whether it was applied. \\
\bottomrule
\end{tabular}
\end{center}

\noindent
To use clean-only forcing, filter on \qcflag{1}. To use the storm-peak velocity moments
recovered from a sensor-block failure, take \qcflag{3} with \code{segValid\_vel} true, knowing
they carry a few-percent scale uncertainty. A segment whose pressure is present but implausible
(inflated depth from a drifting sensor) is flagged \qcflag{4} and contributes to nothing.

At L4 these flags travel per burst as \code{puv\_qc\_flag}, \code{puv\_segValid\_vel}, and
\code{puv\_segValid\_p}; \qcflag{4} segments are dropped from every merged field.

\subsection{Failure signatures}
\label{sec:signatures}

Three distinct failure modes appear in the archive. A one-pass scan of the \code{.sen} files
classifies them without reading the (much larger) \code{.dat} files, because the
\code{.sen} carries battery, sound speed, heading, and tilt:

\begin{center}
\begin{tabular}{@{}l p{7.5cm} c@{}}
\toprule
Signature & How to recognize it & Velocity recoverable? \\
\midrule
Sensor-block failure & Temperature implausible, sound speed steps, battery erratic; tilt stays
  small & \textbf{Yes} (with the rescale) \\
Toppled / moved frame & Temperature, sound speed, battery all normal; tilt grows past
  $30\degree$ or drifts steadily & No \\
Doppler failure & Beam correlation collapses to $\sim$5--10\%, amplitude drops; the auxiliary
  block looks fine & No \\
\bottomrule
\end{tabular}
\end{center}

\noindent
The \code{.sen} scan cannot see a Doppler failure (beam correlation lives in the \code{.dat}),
so ``healthy'' in that scan means the sensor block and frame are healthy, not that the data
are necessarily good.

% ======================================================================
\section{Methods}
\label{sec:methods}

\subsection{L1: quality control and coordinate rotation}

Raw burst files are merged into a continuous 2~Hz series. Clock drift is corrected linearly
over each deployment from the drift measured at recovery. A battery-cutoff detector flags
intermittent sampling by finding gaps over 1~s between consecutive sensor records and
discarding everything after the first such gap. Per-channel QC (\S\ref{sec:qc}) is then
applied. Samples with stable but non-zero tilt (e.g.\ a pipe bent during a storm) are retained
and corrected by a sample-by-sample rotation using the measured pitch $\phi$ and roll $\theta$,
\[
  \mathbf{v}_\text{corrected}
  = \mathbf{R}_\text{roll}(\theta)\,\mathbf{R}_\text{pitch}(\phi)\,\mathbf{v}_\text{raw},
\]
applied before the heading rotation so the velocity is returned to the local vertical frame
before projection onto geographic coordinates. Velocity is then rotated from the instrument
XYZ frame to the buoy frame ($+x$ west, $+y$ north, $+z$ up; left-handed, matching the MOP
convention) using the deployment heading and the IGRF magnetic declination at the deployment
location and epoch. The heading rotation is static (it uses the surveyed heading, not the
live compass), so a failed compass never corrupts it.

\subsection{L2: spectral analysis}

\paragraph{Segmentation.} The L1 series is divided into non-overlapping 1-hour segments (7200
samples at 2~Hz), aligned to the top of the UTC hour. The 1-hour length matches the cadence of
the CDIP MOP wave hindcast and gives the longer averaging window that segment-mean velocity
benefits from (random Doppler error in $\overline{u}$ falls as $1/\sqrt{N}$, so the 7200-sample
noise floor is $\approx 1.9\times$ smaller than a 2048-sample 17-minute window). A stationarity
audit confirmed that bulk parameters are stable on the 1-hour scale at these sites. Segments
with more than 10\% NaN in a channel group are dropped; smaller gaps are linearly interpolated.
Velocity is rotated to a shore-normal frame ($+x$ onshore, $+y$ alongshore north) using the
shore-normal angle from CDIP for the matching MOP station. Depth is
$H = \overline{P}\times 10^4/(\rho g) + z_s$ with $\rho = 1025$~kg\,m$^{-3}$, $g = 9.81$~m\,s$^{-2}$,
and $z_s$ the sensor height above the bed.

\paragraph{Spectral estimation.} Power spectra use the multi-taper method of Thomson (1982)
with time-bandwidth product $NW = 4$ ($K = 7$ Slepian tapers, $\approx 14$ degrees of freedom
per bin). The 1-hour segment gives a native resolution $\Delta f \approx 2.8\times10^{-4}$~Hz.
Pressure--velocity cross-spectra use the same tapers. Multi-taper is chosen for its
bias--variance trade-off: at constant record length it matches Welch's effective bandwidth
while preserving fine frequency structure that Welch averaging smears. (Welch remains available
as a fallback.)

\paragraph{Pressure correction.} The pressure spectrum is converted to surface elevation with
the linear-theory transfer function $K_p(f) = \cosh(kz_s)/\cosh(kH)$, where $k(f)$ solves
$\omega^2 = gk\tanh(kH)$ by Newton's method, and $S_\eta = S_p / K_p^2$. Bins where
$K_p < 0.1$ are zeroed to suppress high-frequency noise amplification; the resulting cutoff
$f_\text{cut}$ is recorded.

\paragraph{Bulk parameters.} Significant wave height is $H_s = 4\sqrt{m_0}$ with
$m_0 = \int S_\eta\,df$, integrated separately over the sea--swell (0.04--0.25~Hz) and
infragravity (0.004--0.04~Hz) bands. Peak period is $T_p = 1/f_p$ from the sea--swell spectral
peak; $T_{m02} = \sqrt{m_0/m_2}$. Mean direction comes from the first-harmonic coefficients
$a_1$, $b_1$ of the pressure--velocity cross-spectra. Energy flux is
$F = \rho g \int c_g(f)\,S_\eta(f)\,df$.

\paragraph{Near-bed velocity, stress, and moments.} Near-bed orbital velocity is obtained by
scaling each velocity-spectrum bin by $\cosh(kH)/\cosh[k(H-z_s)]$ and inverse-transforming;
$U_b$ (rms) and $A_w$ (orbital excursion) follow. Wave bed shear stress uses the Swart (1974)
friction factor with $k_s = 2.5\,D_{50}$ ($D_{50} = 0.25$~mm by default, to be replaced with
site grain size). Reynolds stresses and turbulent kinetic energy come from the detrended
fluctuations. Velocity skewness $S_k$ and asymmetry $A_s$ (Hilbert form) characterize the
wave nonlinearity relevant to sediment transport.

\subsection{L3: wave forcing characterization}

The surface-elevation spectrum is integrated over infragravity (0.004--0.04~Hz), swell
(0.04--0.10~Hz), and sea (0.10--0.25~Hz) bands, with a dominant-band flag by energy flux.
Storm events are found by threshold exceedance ($H_s > 1.5$~m for $\geq$6~h, events within
12~h merged); to characterize forcing during PUV gaps, hourly MOP model data are loaded for a
window extending 7~days past the record, and storms present in MOP but absent from the PUV are
flagged as gap-period storms. Transport proxies include the Shields parameter
$\theta = \tau_b/[(\rho_s-\rho)gD_{50}]$ with the Soulsby--Whitehouse critical threshold and a
mobilization flag, cumulative bottom energy flux, and the Rouse number classifying bedload vs.\
suspended load. Tidal harmonic analysis (\code{t\_tide}, complex $u+iv$) uses the NOAA Scripps
Pier tide prediction as the depth reference ($R = 0.995$ vs.\ $0.987$ for a fit to the noisy
PUV pressure). Subtidal residual currents (observed minus tidal) capture wave-driven mean
flow, including undertow.

% ======================================================================
\section{Validation}
\label{sec:validation}

The pipeline has been validated against the CDIP MOP wave model and against an independent
prior processing pipeline, and its mean-flow output has been tested against physical theory.

\begin{itemize}
\item \textbf{Against MOP.} PUV bulk wave heights agree well with MOP spectra shoaled to the
  instrument depth by linear-theory energy conservation. Residual frequency-dependent biases
  were traced to their sources (nonlinear shoaling, directional narrowing, bound long waves,
  and MOP spectral-peak broadening) rather than left as unexplained scatter.
\item \textbf{Head-to-head cross-validation.} On the legacy Ruby2D cross-shore array, the
  pipeline reproduces an independent prior pipeline's bulk parameters at $R^2 = 0.93$--$0.98$,
  confirming the L1--L2 chain end to end.
\item \textbf{Mean cross-shore flow.} The segment-mean cross-shore velocity was tested across
  38 instruments against three hypotheses --- real Stokes return flow, a constant calibration
  offset, and random Doppler noise. The data match the Stokes return-flow prediction in sign,
  magnitude, and response to averaging-window length, and are incompatible with the offset and
  noise hypotheses. The segment-mean noise floor ($\sim$0.24~mm\,s$^{-1}$ at 1-hour averaging)
  is well below every reported slope and modulation amplitude.
\end{itemize}

\noindent
The deeper validation writeups (\code{results\_validation.tex}, \code{multitaper\_writeup.tex})
carry the full tables, figures, and hypothesis tests.

% ======================================================================
\section{Output reference}
\label{sec:outputs}

Outputs are \code{.mat} files, one per instrument per level, under
\code{outputs/L\{n\}/\{deployment\}/}. The most-used fields:

\paragraph{L1 (\code{PUV} struct).} \code{time} (naive-UTC datetime), \code{fs}, \code{P}
(dBar), \code{BuoyCoord.\{U,V,W\}} (geographic velocity), \code{T}, \code{doffp}, and the
\code{PUV.qc} provenance sub-struct (\code{valid\_vel}, \code{valid\_p}, \code{valid\_tilt},
\code{valid\_T}, \code{vel\_c\_factor}, \code{vel\_c\_corrected}, \code{vel\_rotation\_static}).

\paragraph{L2 (\code{L2} struct, one row per segment).} \code{time}, \code{Hs},
\code{Hs\_SS}, \code{Hs\_IG}, \code{Tp}, \code{Tm02}, \code{meanDir}, \code{Ef}, \code{depth},
\code{Ub}, \code{tau\_b}, the \code{vmom} sub-struct (\code{skewness}, \code{asymmetry},
\code{u\_uabs2}, \dots), \code{uMean}/\code{vMean}, and the provenance fields of
\S\ref{sec:qc} (\code{segValid}, \code{segValid\_vel}, \code{segValid\_p}, \code{qc\_flag},
\code{Hs\_source}).

\paragraph{Coordinate conventions.} After L1 rotation: $+x$ west, $+y$ north, $+z$ up. After
L2 shore-normal rotation: $+x$ onshore, $+y$ alongshore north; onshore flux is $q_x > 0$.
Pressure is stored in dBar.

% ======================================================================
\section{Troubleshooting and where to look next}
\label{sec:troubleshooting}

\begin{itemize}
\item \textbf{A driver errors that the previous level is missing.} Run L1$\rightarrow$L4 in
  order; each reads the level below it.
\item \textbf{Network reads are slow.} Cache raw files with \code{copy\_raw\_to\_local(cfg)};
  the drivers use \code{raw\_cache/} automatically.
\item \textbf{A deployment recovers little velocity during a storm.} Check the failure
  signature (\S\ref{sec:signatures}) from the \code{.sen}: a sensor-block failure is
  recoverable, a topple is not. Confirm \code{cfg.qcOpts.Tvalid} is set to the site range.
\item \textbf{Recovered ($\qcflag{3}$) data in an analysis.} Decide deliberately: filter to
  \qcflag{1} for clean-only, or include \qcflag{3} knowing the scale uncertainty.
\end{itemize}

\noindent
Deeper references:
\code{docs/pipeline\_levels.md} (level-by-level detail and output-struct shapes),
\code{PIPELINE\_NOTES.md} (design decisions, coordinate conventions, known issues),
\code{docs/OUTSTANDING\_channel\_decoupling.md} and
\code{docs/L1\_sensor\_block\_failure\_2026-07-09.md} (the channel-decoupling work in full),
\code{config/CONFIG\_REVIEW\_NOTES.md} (per-deployment review items).

\vfill
\noindent\rule{\textwidth}{0.4pt}\\[2pt]
\small\emph{Maintainer: Holden Leslie-Bole, Scripps Institution of Oceanography.}

\end{document}
